% HD-G25. GPT-6 Astra Pro, /root/g34; 21 November 2026, inherited effort.
\section{HD-G25: the exhaustive natural Gaussian reflection and scalar source filter}
\label{sec:gaussian-reflection-classification}

\paragraph{The actual source class.}
Put \(K=\mathbb Q(i)\). Let \(W\) be a finite irreducible
reflection group in its natural representation, and let
\[
             G=\langle W,Z\rangle\ \subset GL_r(K)
\]
be an actual matrix group, where \(Z\) consists of finite-order
scalar matrices in \(K\). Use the natural primitive rational
group-algebra factor, with corner \(K\), of the Jacobian of a
connected full-\(G\) sphere cover. Its Gaussian action is the
action from this specified corner, with the inverse deck convention
already fixed in HD-S23. Suppose the factor is a Weil sixfold
of signature \((3,3)\), with the restricted theta polarization.

Every rank-nine period family from this source class has ordinary
Hermitian determinant \([-1]\) or \([-3]\). This exhausts the
natural Gaussian reflection representations and their
\(K\)-defined scalar extensions in the stated primitive-corner
sense. The proof includes the exact coefficient changes and
polarized maps. It makes no assertion about nonscalar overgroups
or an additional Gaussian action imposed only on a special
subfamily of a real primitive factor.

\paragraph{1. The published classification and its exact Gaussian restriction.}
Cohen, \emph{Finite complex reflection groups}, Ann.\ Sci.\ ENS
9 (1976), gives the imprimitive classification in Theorem (2.4),
printed pp.386--387; the real Coxeter list in (4.4)(i), p.399;
and the nonreal primitive classification in Theorem (5.12),
p.435, with its complete list in table (4.16), p.412.
His word \emph{complex} in the last theorem means nonreal, as
defined in (1.1), p.380. The primary source is
\[
 \text{\url{https://www.numdam.org/item/ASENS_1976_4_9_3_379_0.pdf}}.
\]
The cited original pages and the relevant root pages were
reopened and checked. We now give the exact trace tests that
restrict this classification to a \(K\)-matrix model.

For an imprimitive natural representation of rank \(r\ge3\),
the classification gives \(G(m,p,r)\), \(p\mid m\), with
the usual monomial matrices. Write
\(\zeta=\exp(2\pi i/m)\). Its subgroup \(G(m,m,r)\)
contains the two literal diagonal matrices
\[
 \operatorname{diag}(\zeta,\zeta^{-1},1,\ldots,1),\qquad
 \operatorname{diag}(\zeta,\zeta,\zeta^{-2},1,\ldots,1).
\]
The trace of the first makes
\(\zeta+\zeta^{-1}\in K\cap\mathbb R=\mathbb Q\).
It is an algebraic integer in \([-2,2]\), hence one of
\(-2,-1,0,1,2\). Consequently
\(m\in\{1,2,3,4,6\}\). The case \(m=1\) is reducible
in the natural rank-\(r\) permutation representation.
For \(m=3\) the second trace is \(r-3+3\zeta\);
for \(m=6\) it is \(r-3+\zeta\).
Neither lies in \(K\). Thus only
\[
                           m=2,\ 4                         \tag{G25.1}
\]
remain. This argument uses traces of actual group elements,
which are unchanged by a complex change of basis.

For the nonreal primitive list, the exact exclusions are:
\[
\begin{array}{c|l}
\text{Shephard--Todd group}&\text{trace of an actual element}\\ \hline
G_{24}&(-1+\sqrt{-7})/2\\
G_{25},G_{26}&2+\omega\ \text{or }2+\bar\omega\\
G_{27}&(1+\sqrt5)/2\\
G_{32}&3+\omega\ \text{or }3+\bar\omega\\
G_{33}&2+3\omega\\
G_{34}&6\omega
\end{array}                                                \tag{G25.2}
\]
Here \(\omega^2+\omega+1=0\).
Every listed number is outside \(K\); for the real entries
this follows from \(K\cap\mathbb R=\mathbb Q\), and the
two imaginary quadratic fields have intersection \(\mathbb Q\)
with \(K\).

The trace assertions have the following concrete sources and
maps. For \(G_{24}\), the exact matrices and the intertwiner
to Cohen's roots in HD-S13 give
\(\operatorname{tr}((s_1s_2s_3)^2)=(-1+\sqrt{-7})/2\).
The \(G_{25},G_{26},G_{32}\) entries are traces of the
order-three reflections displayed in Cohen's root diagrams and
table: their eigenvalues are \(1\) in all but one coordinate
and a primitive cube root in the remaining coordinate.
For \(G_{27}\), the two roots on the labelled five-edge in
the \(J_3(5)\) diagram, p.408, have inner-product absolute
value \(\cos(\pi/5)\). The product of their two involutive
reflections has two eigenvalues \(e^{2\pi i/5},e^{-2\pi i/5}\)
and fixes the orthogonal complement. Its trace is the value
in (G25.2).

For \(G_{33}\), Cohen (5.10), printed p.430, explicitly
embeds \(G(3,3,4)\) as a four-dimensional reflection subgroup.
The element
\(\operatorname{diag}(\omega,\omega,\omega,1)\) of that
subgroup fixes the remaining fifth coordinate; its trace is
\(2+3\omega\). Finally Cohen (1.4)(ii), p.381, and the
degrees of \(G_{34}\) in table (4.16) give scalar centre
of order six. It contains the literal scalar \(\omega I_6\),
with trace \(6\omega\).

The real Coxeter groups of rank at least three are
\(A_r,B_r,D_r,E_6,E_7,E_8,F_4,H_3,H_4\), with the
usual small-rank coincidences allowed. The five-edge in
\(H_3,H_4\) gives the same two-reflection product calculation:
its trace is
\[
                  r-2+2\cos(2\pi/5)
                    =r-\frac52+\frac{\sqrt5}{2},
\]
which is outside \(K\). All other real groups in this list
have their usual rational root representations. The only
nonreal primitive Gaussian cases left by (G25.2) are therefore
\(G_{29}\) and \(G_{31}\), already given by literal
Gaussian matrices in HD-S32.

\paragraph{2. The exact \(K\)-intertwiner and the scalar extension.}
The classification initially supplies a complex conjugacy.
For every retained case the preceding paragraphs and HD-S32
give an actual standard model \(\rho_0\) over \(K\).
Let \(\rho\) denote the supplied \(K\)-model, with the group
isomorphism to \(\rho_0\) fixed by that classification conjugacy.
Form the finite \(K\)-linear system
\[
                   T\rho(g)=\rho_0(g)T
                         \quad(g\text{ in a generating set}).
                                                            \tag{G25.3}
\]
Its complex solution space is one-dimensional and contains
an invertible conjugacy. Gaussian elimination over \(K\)
has the same rank as over \(\mathbb C\). Choose its first
nonzero free-coordinate solution by setting that coordinate
to one and the other free coordinates to zero. This specifies
an actual nonzero \(K\)-matrix \(T\). Any nonzero solution
is invertible: its kernel and image are invariant, and both
representations are irreducible. Equation (G25.3) is the
required \(K\)-linear intertwiner.

This construction can be made without knowing the entries
of a complex conjugacy: enumerate the finitely many bijections
between the two finite groups, retain those preserving
multiplication, and solve (G25.3) for each. The classification
guarantees an invertible solution for at least one retained
bijection. Taking the first such solution is a finite
procedure entirely over \(K\).

Every finite-order scalar in \(K\) belongs to \(\mu_4\).
Indeed it is an integral Gaussian number of norm one.
The same \(T\) intertwines these literal scalar matrices,
so the scalar extension is also transported by (G25.3).
In particular the standard monomial coordinate lines in
(G25.1) pull back through \(T\) to an actual \(K\)-defined
line system, exactly the hypothesis of HD-S34.

For a real retained group without the scalar \(i\), the
rational matrix model has simple rational block
\(\operatorname{Mat}_r(\mathbb Q)\). One way to check the
map is to apply the rational representation of the group
algebra: after scalar extension its image is the full complex
matrix algebra by absolute irreducibility, so its rational
image already has dimension \(r^2\). Thus its primitive
corner is \(\mathbb Q\). Adjoining scalar \(-1\) does
not change this image. No element of that corner squares
to \(-1\), so it is not the specified Gaussian corner.
The real cases relevant to the assertion consequently use
the actual scalar extension by \(i\).

The coefficient-form comparison is also explicit. Pull a
standard positive invariant \(K\)-Hermitian form back along
\(T\). Relative to any other such form its Gram operator
commutes with \(\rho(G)\); absolute irreducibility makes
it a scalar \(a\in K\cap\mathbb R=\mathbb Q_{>0}\).
Apply \(T\) to every coordinate in the parabolic cochain
complex. This gives an actual isomorphism of paired
parabolic spaces, up to the scalar \(a\). In dimension six
its determinant multiplier is \(a^6=N_{K/\mathbb Q}(a^3)\).
Changes of \(K\)-basis contribute the norm of their determinant.

The coefficient-to-Jacobian maps of HD-S23 compose with
this cochain map, with their exact Schur factors retained.
For a chosen coefficient vector \(v_\nu\) of positive rational
norm \(q_\nu\), use
\[
 j^\nu_x(g)=h_\nu(v_\nu,\rho_\nu(g)x)/q_\nu,\qquad
 \Phi_\nu=\operatorname{PD}(j^\nu+\bar j^\nu),\qquad
 c_\nu=\frac{2|G|}{r q_\nu}.
\]
The two conjugate cross terms give
\(h_{B_\nu}(\Phi_\nu x,\Phi_\nu y)=c_\nu H_\nu(x,y)\).
Choose a positive integer \(M\) carrying the source homology
lattice into the target under
\[
                  F=M\Phi_0T\Phi_{\rm given}^{-1}.
\]
This explicit Hodge map defines the homomorphism on the
universal covers of the corresponding complex tori. It is
a rational homology isomorphism, hence an isogeny of the
abelian factors. Its literal polarization scalar is
\[
                    M^2a\,c_0/c_{\rm given},                \tag{G25.3a}
\]
which is positive and rational; its sixth power is a norm.
Thus both the coefficient scalar and the two Schur factors
are accounted for, including the integer clearing the lattices.
These fixed maps also intertwine the period maps. The
determinant and rank conclusions below are transported by
these specified coefficient and abelian-variety maps.

\paragraph{3. The additional \(F_4\) case.}
This case admits a stronger rank-nine conclusion than a
group-order denominator estimate. In the standard rational
space with Gram \(I_4\), take the four simple roots
\[
 e_2-e_3,\quad e_3-e_4,\quad e_4,\quad
                \tfrac12(e_1-e_2-e_3-e_4).
\]
Their Gram matrix is
\[
 \begin{pmatrix}
 2&-1&0&0\\-1&2&-1&0\\0&-1&1&-1/2\\0&0&-1/2&1
 \end{pmatrix}.
                                                            \tag{G25.4}
\]
The unit-root graph is precisely the \(3,4,3\) Coxeter
graph \(F_4\) on Cohen's p.399. The actual reflections are
\(x\mapsto x-2(a^tx)/(a^ta)a\).
They have entries in \(\tfrac12\mathbb Z\), preserve \(I_4\),
and preserve the explicit 48-root set
\[
 \{\pm e_j\}\ \cup\
 \{\pm e_j\pm e_k:j<k\}\ \cup\
 \{\tfrac12(\pm e_1\pm e_2\pm e_3\pm e_4)\}.
\]
The first two reflections interchange coordinates \(2,3,4\),
and the third changes the sign of \(e_4\). Thus the sign
changes in coordinates \(2,3,4\) lie in the generated group.
If \(r_4\) is the fourth reflection and \(s_j\) the sign
change in coordinate \(j\), the literal vector calculation
\[
                  r_4s_2s_3r_4(e_4)=e_1
\]
conjugates the third reflection to the sign change in
coordinate \(1\). Their product is \(-I_4\).
Adjoining \(iI_4\) gives
\[
                  G_F=\mu_4W(F_4),\qquad G_F/W(F_4)\simeq C_2.
\]
These are literal rational and Gaussian matrix models.
The exact root-closure certificate independently gives
\(|W(F_4)|=1152\), \(|G_F|=2304\), absolute character
norm one and entries in \(\tfrac12\mathbb Z[i]\).
The universal argument below uses the displayed roots and
matrices, not a sample of branch tuples.

Let \(g_1,\ldots,g_b\) be a true-branch generating tuple in
\(G_F\) giving the natural sixfold. Its exact parabolic rank
formula is
\[
              \sum_{j=1}^b\operatorname{rank}(1-g_j)=14.
                                                            \tag{G25.5}
\]
Each summand is at least one. A rank-nine period map requires
\(b\ge12\), since the Hurwitz configuration dimension is \(b-3\).

There are at least two entries outside \(W(F_4)\): the
tuple generates the quotient \(C_2\), and its product is
one. Such an entry is \(g=iw\) with \(w\) real orthogonal.
Its fixed space is the \(-i\) eigenspace of \(w\); complex
conjugation pairs this with the \(i\) eigenspace of the
same dimension. Consequently its fixed dimension is at
most two, and its codimension is at least two. The two
such entries already use all of the excess
\(\sum(\operatorname{codim}-1)=14-b\le2\).
Therefore \(b=12\), exactly two entries have odd scalar
phase, each has codimension two, and all other entries
are real reflections.

For either odd entry the two eigenspaces just described
have dimensions two and two and exhaust the four-dimensional
space. Thus \(w^2=-I_4\) and \(g^2=I_4\).
Every other entry is an involutive real reflection.
Hence every local inertia is an involution.

Now use \(R=\mathbb Z[i,1/2]\), \(L=R^4\) and Gram
\(I_4\). The actual local projector is \((1+g_j)/2\),
and \((1-g_j)/2\) is its complementary image. The central
scalar \(i\), present in full monodromy, kills absolute
coinvariants because \(i-1\) is a unit in \(R\), for both
the coefficient system and its dual. The complete primitive-image
cap-product argument of HD-S27 therefore gives a perfect
rank-six parabolic skew-Hermitian form over \(R\).
Its Hermitian form \(iS\) has negative rational determinant
in \(R^\times\), hence determinant \(-2^a\).
The identity \(2=N(1+i)\) gives ordinary class \([-1]\).

For completeness the actual primitive-factor scalar is
fixed as follows. Take \(v=(1+i)e_1\), of norm two,
and
\[
 e=\frac4{|G_F|}\sum_{g\in G_F}
       \operatorname{Re}h_0(v,\rho(g^{-1})v)\,g,
 \qquad D=|G_F|/4=576.
\]
The half-Gaussian entry bound gives \(De\in\mathbb Z[G_F]\).
Matrix-coefficient orthogonality makes \(e\) primitive and
self-adjoint with corner \(K\). On
\(B=\operatorname{im}(De)\subset J(C)\), use inverse
scalar deck multiplication by \(i\).
The literal map
\[
 j_x(g)=h_0(v,\rho(g)x)/2,\qquad
                \Phi=\operatorname{PD}(j+\bar j)
\]
has complex Schur scalar \(|G_F|/8\); its two conjugate
cup terms give
\(h_B(\Phi x,\Phi y)=D(iS)(x,y)\), with exactly the
left-deck convention of HD-S23. The multiplier
\(D^6=N(D^3)\) preserves the determinant class.
Thus every rank-nine \(F_4\) scalar source in the assertion
is split, for its actual restricted polarization.

\paragraph{4. Exhausting the retained cases.}
For coefficient rank one or two, the sphere parabolic formula
is \(\sum c_j=2r+6\). Faithfulness and true branching give
\(c_j\ge1\), so the period ranks are at most five and seven,
respectively. These bounds use the actual cover-to-configuration
and relative-Jacobian maps, as in HD-S30.

For rank at least three the complete list and the corresponding
proved source conclusions are
\[
\begin{array}{c|l|l}
\text{natural reflection type}&\text{actual source result}
                    &\text{rank-nine determinant}\\ \hline
G(m,p,r),\ m=2,4 &\text{HD-S34, \(K\)-defined monomial lines}&[-1]\\
A_r &\text{HD-S33, scalar augmentation of \(S_{r+1}\)}&[-1]\\
E_6,E_7 &\text{HD-S27, exact integral parabolic maps}&[-1]\\
E_8 &\text{HD-S28, every profile with \(b\ge12\)}&[-1]\text{ or }[-3]\\
F_4 &\text{paragraph 3 above}&[-1]\\
G_{29},G_{31}&\text{HD-S32, every branch profile}&[-1]\text{ or }[-3]
\end{array}                                                  \tag{G25.6}
\]
The real types \(B_r,D_r\) are already in the first row.
For \(A_r\), \(r\ge2\), the natural augmentation permutation
group contains no nonidentity scalar: a scalar \(-1\)
would have impossible permutation trace \(2-(r+1)<0\)
on the full permutation module. Thus adjoining \(\mu_4\)
is literally the direct product used in HD-S33, rather
than an unidentified central quotient.

The standard monomial lines in the first row are genuinely
defined over \(K\), by (G25.3). HD-S34 constructs their
coset basis, local cycle maps and global holonomy relation,
the primitive subgroup idempotent, an integral Gaussian
action from a stabilizer element, and the actual cyclic
quotient/polarized pullback. Its conclusion therefore
applies with no extension of the coefficient field.
All other rows likewise retain their actual coefficient
and primitive-Jacobian maps. The ordinary split-class
identification with \(U^3\) is the proved isometry
construction of HD-S03a.

This proves the stated exhaustive determinant filter for
rank-nine families. Compatible flat-polarization changes
under a specified fixed isogeny preserve it by HD-S17 on
the open unitary period domain: their self-adjoint scalar
is rational and positive, with a norm sixth power.
No arbitrary polarization change at an isolated special
point is included.

\paragraph{Check and source scope.}
The classification inputs are the published complete
statements and the explicit trace obstructions above.
The \(F_4\) certificate \path{f4_certificate.py}
constructs its 48 roots and four reflection matrices from
scratch, enumerates their faithful root permutations,
then adjoins the scalar on the 96 real and imaginary
roots. It checks the two orders, the actual central
scalar, half-integral entries and character norm.
It is an exact finite algorithm, not a proof-assistant
certificate and not a search over unlisted reflection groups.
The other source rows refer to their complete, separately
checked proofs. The result supplies a closed source filter;
it adds no cycle, dominance or all-finite-group assertion.
